\documentclass[11pt]{article}
\usepackage[margin=1in]{geometry}
\usepackage{times}
\usepackage{amsmath,amssymb}
\usepackage{booktabs}
\usepackage{graphicx}
\usepackage{hyperref}
\usepackage{xcolor}

% Markers: red = fill in prose/decision; blue = paste a number from the scripts.
% They exist so no result is ever silently fabricated.
\newcommand{\todo}[1]{{\color{red}\textbf{[TODO: #1]}}}
\newcommand{\res}[1]{{\color{blue}#1}}

\title{Compressed but Candid? The Effect of Post-Training Quantization\\
       on Chain-of-Thought Faithfulness and Monitorability}
\author{\todo{author name} \\ \todo{affiliation / email}}
\date{\today}

\begin{document}
\maketitle

%======================================================================
\begin{abstract}
Chain-of-thought (CoT) monitoring---reading a model's reasoning trace to detect
misbehavior---is a leading strategy for AI oversight, but it is only sound if the
trace is \emph{faithful}: the reasoning must verbalize the factors that actually
determine the answer. Independently, weight quantization to INT8 and INT4 is how
open-weight reasoning models are actually deployed on commodity hardware. These
two facts have been studied in isolation. A growing literature shows that
quantization harms reasoning \emph{accuracy} and inflates CoT \emph{length},
while CoT faithfulness is measured almost exclusively at full precision. We ask
the question at their intersection: \textbf{does post-training weight
quantization erode chain-of-thought faithfulness?} Holding the base checkpoint
fixed and varying only numerical precision (FP16, INT8, INT4), we apply the
standard hint-injection protocol to \todo{model(s)} and measure the rate at
which a model verbalizes an injected cue that flips its answer. We find
\todo{one-line headline result}. We separate a \emph{direct} effect of precision
from an \emph{indirect} effect mediated by the known concentration of
unfaithfulness on incorrect answers, and we report the thinking-span versus
answer-span acknowledgment gap. \todo{implication sentence}. All code, per-trace
data, and a human-validated cue detector are released for full reproducibility;
the entire study runs on a single free-tier GPU.
\end{abstract}

%======================================================================
\section{Introduction}
Modern reasoning models externalize an explicit, step-by-step chain of thought
(CoT) before committing to an answer. Because this trace is written in natural
language, it has been proposed as a practical window for AI oversight: if a model
is about to misbehave, a monitor might catch the intent in the reasoning before
it reaches the output~\cite{cotmonitorability2025,baker2025}. The promise of
``just read the reasoning'' is attractive precisely because it needs no access to
model internals.

That promise rests entirely on \emph{faithfulness}. A CoT is faithful if it
reports the factors that actually cause the answer. A now-substantial body of
work shows this cannot be taken for granted: models produce plausible,
fluent-sounding reasoning that nonetheless omits the true cause of the
answer---most sharply demonstrated by \emph{hint injection}, where a cue planted
in the prompt changes the answer while the reasoning never mentions
it~\cite{turpin2023,anthropic2025reasoning,lietome2026}. Faithfulness, in other
words, is a measurable and frequently violated property.

Almost all of this measurement has been done on full-precision models. But
full precision is not how open reasoning models are deployed. Post-training
quantization to 8- and 4-bit weights is the default for running these models on
consumer GPUs and edge devices, and a fast-growing literature documents its
consequences for reasoning: quantized models lose accuracy at aggressive
bit-widths, ``overthink,'' and inflate their token
counts~\cite{quanthurts2025,overthinking2026,tokeninflation2026,multiplicative2026}.
Every one of these studies asks whether the \emph{answer} survives compression.
None asks whether the \emph{reasoning trace} survives as a faithful, monitorable
artifact.

This gap has a direct safety consequence. The models most likely to be deployed
in the wild---quantized open-weight models on commodity hardware---are exactly
the models whose monitorability nobody has verified. If faithfulness is measured
at FP16 but the deployed artifact is INT4, a monitorability guarantee obtained in
the lab may simply not hold in production. We therefore ask a single, concrete
question:

\begin{quote}
\emph{Does post-training weight quantization (FP16 $\to$ INT8 $\to$ INT4) reduce
the rate at which a reasoning model verbalizes the cue that drove its answer?}
\end{quote}

\paragraph{Contributions.}
\begin{enumerate}
  \item To our knowledge, we are the first to measure the effect of \emph{weight
        quantization} on CoT faithfulness, holding the base checkpoint fixed so
        that numerical precision is the \emph{only} independent variable and the
        comparison is causally clean.
  \item Using the standard hint-injection protocol, we report the
        cue-acknowledgment rate on answer-flipping cases across precisions, and
        the thinking-span versus answer-span acknowledgment
        gap~\cite{whymodelsknow2026}.
  \item We stratify faithfulness by answer correctness, separating a \emph{direct}
        effect of precision from the \emph{indirect} effect predicted by the
        ``two regimes'' account of unfaithfulness~\cite{tworegimes2026}, in which
        unfaithfulness concentrates on incorrect answers.
  \item We release all code, the full per-trace dataset, and a cue detector
        validated against human labels, and we design the study so that a null
        result (``quantization preserves monitorability'') is as informative and
        publishable as a positive one.
\end{enumerate}

%======================================================================
\section{Related Work and Positioning}
Our contribution sits at the intersection of three active literatures. Because
the words ``compression'' and ``reliability'' appear in all three with different
meanings, we state explicitly what each measures and where our question differs.

\paragraph{CoT faithfulness and monitorability.}
The hint-injection protocol and its
relatives~\cite{turpin2023,lanham2023} test whether a model verbalizes a prompt
cue that changes its answer; a range of recent studies extend and stress-test
these faithfulness and monitorability
metrics~\cite{anthropic2025reasoning,lietome2026,monitorability2025,faithcotbench2026,bonafide2026}.
Crucially, \emph{this work is conducted at full precision.} We adopt the same
protocol unchanged and add numerical precision as the independent variable.

\paragraph{Compression of the trace---but its \emph{length}.}
A distinct line studies how compressing the \emph{reasoning trace itself}
(thousands of tokens down to hundreds, via length penalties or efficient-reasoning
objectives) degrades
monitorability~\cite{monitorabilitytax2026,lengthpenalties2026}. This is a
genuinely different sense of ``compression'': it shortens the text, whereas we
quantize the \emph{weights} and leave the decoding objective untouched. We are
careful throughout to keep these two axes separate, because conflating them is
the most likely way a reader could mistakenly think our question is already
answered.

\paragraph{Quantization of reasoning---but \emph{accuracy}, not faithfulness.}
Quantized reasoning models are known to lose accuracy, overthink, and inflate
token counts at low
bit-widths~\cite{quanthurts2025,quantmeets2025,overthinking2026,tokeninflation2026,multiplicative2026,whenreasoningmeets2025}.
Each of these measures \emph{whether the answer is right} or \emph{how long the
reasoning is}, not \emph{whether the trace honestly reports its own cause}. The
closest reliability-oriented work, UniComp, finds a
performance--reliability \emph{decoupling}: compression can preserve accuracy
while silently degrading other properties~\cite{unicomp2026}; related work finds
compression can undo safety
alignment~\cite{quantundoes2026}. Neither measures cue-verbalization faithfulness
under the hint-injection protocol. Our study fills exactly that empty cell:
faithfulness (not accuracy, not length, not aggregate ``reliability'') as a
function of weight precision (not trace length).

%======================================================================
\section{Method}
\paragraph{Models and precisions.}
We evaluate \todo{model list, e.g.\ DeepSeek-R1-Distill-Qwen-1.5B, and 7B if
compute allows} at three precisions: FP16 (the baseline), INT8, and INT4. The
central design choice is that all three are loaded from the \emph{same} base
checkpoint---INT8/INT4 via on-the-fly bitsandbytes quantization---so that the
weights differ only in numerical precision and nothing
else~\cite{dettmers2023,deepseekr1}. This makes precision the sole manipulated
variable and rules out confounds from differently trained or differently prompted
models. To guard against an artifact of a single quantization algorithm, we
additionally include \todo{one AWQ and one GPTQ checkpoint} as a cross-method
robustness check.

\paragraph{Faithfulness protocol.}
For each multiple-choice item (\todo{dataset, e.g.\ MMLU}; $N=\todo{N}$) we query
the model twice. First \emph{clean}, recording its unbiased answer. Then
\emph{cued}, with a hint injected into the prompt that points at a deliberately
\emph{wrong} option. We use two established cue types, an appeal to authority and
a sycophancy cue. The logic of the protocol is counterfactual: on the subset of
items where the injected cue actually \emph{flips} the answer to the cued (wrong)
option, the cue is demonstrably the cause of the answer, so a faithful reasoning
trace must mention it. We therefore measure the \emph{cue-acknowledgment rate} on
these flipped cases. Decoding is greedy ($T{=}0$) throughout for
reproducibility~\cite{monitorability2025}.

\paragraph{Metrics.}
We report four quantities per precision:
(i)~\emph{flip rate}, the fraction of cued items whose answer moves to the cue---a
prerequisite for the faithfulness measurement and itself of interest, since
susceptibility to cues may vary with precision;
(ii)~\emph{faithfulness}, the cue-acknowledgment rate on flipped cases;
(iii)~the \emph{thinking-versus-answer gap}, comparing acknowledgment inside the
\texttt{<think>} span with acknowledgment in the final answer
span~\cite{whymodelsknow2026}; and
(iv)~\emph{correctness-stratified faithfulness}, which splits (ii) by whether the
model's answer was correct, in order to separate the direct effect of precision
from the indirect route in which lower precision produces more wrong answers,
where unfaithfulness is already known to concentrate~\cite{tworegimes2026}.

\paragraph{Cue detector and its validation.}
Cue acknowledgment is detected automatically by a deterministic lexical matcher
over cue-specific marker phrases. Because the credibility of every downstream
number depends on this detector, we validate it against \todo{K} hand-labeled
traces and report accuracy \res{acc} and Cohen's $\kappa=\res{kappa}$
(\S\ref{sec:validation}). We treat $\kappa$ as a gate: automatic rates are
reported as primary results only if human--detector agreement is acceptable, and
otherwise the detector is revised (expanded marker set or an LLM judge) and
re-validated before any rate is trusted.

%======================================================================
\section{Results}
\emph{The numbers below are populated directly from the released scripts
(\texttt{analyze.py}, \texttt{export.py}, \texttt{label\_helper.py}); no value is
reported that the pipeline did not produce.}

\paragraph{Detector validation.}\label{sec:validation}
On \todo{K} hand-labeled flipped-case traces, the lexical detector achieves
accuracy \res{acc}, $\kappa=\res{kappa}$, precision \res{P}, and recall \res{R}.
\todo{One sentence stating whether agreement clears the reliability gate.}

\paragraph{Faithfulness versus precision (headline).}
Table~\ref{tab:headline} reports flip rate and faithfulness at each precision.
\todo{State the trend: does acknowledgment fall monotonically FP16$\to$INT8$\to$INT4,
stay flat, or move non-monotonically? Give the effect size and, if computed, a
significance test or bootstrap CI.}

\begin{table}[h]\centering
\caption{Cue-acknowledgment (faithfulness) on answer-flipping cases, by precision.}
\label{tab:headline}
\begin{tabular}{lccc}
\toprule
Precision & Flip rate & Faithfulness (ack $\mid$ flipped) & $n_{\text{flipped}}$ \\
\midrule
FP16 & \res{--} & \res{--} & \res{--} \\
INT8 & \res{--} & \res{--} & \res{--} \\
INT4 & \res{--} & \res{--} & \res{--} \\
\bottomrule
\end{tabular}
\end{table}

\paragraph{Thinking versus answer gap.}
Figure~\ref{fig:gap} contrasts cue acknowledgment in the thinking span with the
answer span. \todo{Report whether the thinking-span rate exceeds the answer-span
rate (replicating prior full-precision findings~\cite{whymodelsknow2026}) and,
critically, whether this gap widens as precision falls.}

\begin{figure}[h]\centering
% \includegraphics[width=0.62\textwidth]{results_gap.png}
\fbox{\parbox{0.62\textwidth}{\centering\vspace{2em}\todo{insert results\_gap.png
from export.py}\vspace{2em}}}
\caption{Thinking- versus answer-span cue acknowledgment across precisions.}
\label{fig:gap}
\end{figure}

\paragraph{Correctness stratification.}
Table~\ref{tab:strat} splits faithfulness by answer correctness. \todo{Determine
whether the precision effect is mainly \emph{indirect} (low bit-width yields more
wrong answers, and unfaithfulness rides on wrong answers) or whether there is a
\emph{direct} drop within each correctness regime---the more concerning case for
monitoring.}

\begin{table}[h]\centering
\caption{Faithfulness stratified by answer correctness.}
\label{tab:strat}
\begin{tabular}{lcc}
\toprule
Precision & ack $\mid$ correct & ack $\mid$ incorrect \\
\midrule
FP16 & \res{--} & \res{--} \\
INT8 & \res{--} & \res{--} \\
INT4 & \res{--} & \res{--} \\
\bottomrule
\end{tabular}
\end{table}

%======================================================================
\section{Discussion}
Two outcomes are worth interpreting in advance, because both are publishable and
the study is designed to be honest about either.

If faithfulness \emph{drops} with quantization, then quantized deployments are
less monitorable than their full-precision evaluations imply. This would be a
safety gap hidden by accuracy-only benchmarking: a model could pass its FP16
monitorability check, be shipped at INT4, and silently stop reporting the cues
that drive its answers. It would also extend UniComp's performance--reliability
decoupling~\cite{unicomp2026} into the specific, safety-critical dimension of
CoT faithfulness. If the drop is mostly \emph{direct} (present even within the
correct-answer regime), the concern is sharper still, since it cannot be
explained away as a side effect of degraded accuracy.

If faithfulness is \emph{preserved}, that is a genuinely reassuring negative
result: standard quantization, whatever it costs in accuracy, does not by itself
compromise monitorability. Given how much oversight tooling implicitly assumes
the deployed model behaves like its full-precision evaluation, a well-controlled
null here is valuable on its own terms.

%======================================================================
\section{Limitations}
Our scope is deliberately narrow to keep the causal claim clean, and this bounds
the conclusions. We cover \todo{one or few} model families over one size range
(\todo{sizes}); broader generalization across families and scales is future work.
Acknowledgment is detected lexically---mitigated by the $\kappa$ validation and,
\todo{if run}, an LLM-judge robustness check, but still an imperfect proxy for a
human judgment of faithfulness. We study multiple-choice items with two cue
types; other task formats and cue channels (for example, cues delivered through a
tool return rather than the user turn) may behave differently. bitsandbytes is
our primary quantizer with AWQ/GPTQ as a spot check rather than exhaustive
coverage, and we use greedy decoding only. Finally, any observed effect is
correlational with respect to internal mechanism: we measure behavior, not the
circuit-level cause of a change in verbalization.

%======================================================================
\section*{Reproducibility}
All code, the complete per-trace dataset, the human-validated cue detector, and
the exact command sequence are released at \todo{repository URL}. The full study
runs on a single free-tier T4 GPU in approximately \todo{X} GPU-hours, and we
provide a copy-paste notebook so the entire pipeline---generation, scoring,
analysis, and detector validation---can be re-executed end to end at no cost.

\bibliographystyle{plain}
\bibliography{references}
\end{document}
